\subsection{Okubo-Weiss Corner Branch Selection (Deferred 2026-07-02)}
\label{sec:ch9:ow_corner_branch}

% Research note carried over unchanged from phd_thesis_template.tex's
% prior Chapter 5 (git history, this file). Relocated here because it
% is a specific, evaluable future-work item, not a general limitation.
% Sources: experiments/ow_clean_runs.py and
% experiments/outputs/ow_clean/ under
% trunk/Python_Simulations/Vector_Fields/VF_Robot/, git commits 7dc01a8
% and c19732f, and plan.md (repo root) Phase 3 and Phase 7 entries.

\textbf{Context.} The Okubo-Weiss boundary tracker of the second paper
drives $\hat{D} = \det(\hat{\mathbf{J}})$ to zero and drifts along the
level-set tangent $\mathrm{perp}(\nabla\hat{D})$, sign-aligned with the
ambient flow. On the steady double gyre the $D = 0$ set is not a pair of
closed diamonds but a network of straight lines,
$x_f \pm y_f \in \tfrac{1}{2} + \mathbb{Z}$, which cross at the diamond
corners where $\nabla D = \mathbf{0}$. Noise-free experiments (four
starts) showed lock-on within 13 to 16 control steps and mean $|D|$
between 0.003 and 0.007 on the boundary, 0.3 to 0.8 percent of the well
depth $\pi^4 A^2 = 0.974$, but the tangent follows its line straight
through every corner crossing: the gradient-floor fallback carries the
team across the degenerate point and the Newton action resumes on the
continuing line. Circulation around a single vortex core therefore never
occurs with the published law. The paper presents this straight-through
behavior honestly; the wording of Problem 2 was changed from ``circulate
along'' to ``travel along'' the boundary.

\textbf{The deferred idea.} At an X-crossing the field is locally
$D(\mathbf{p}^* + \mathbf{t}) \approx \tfrac{1}{2}\mathbf{t}^\top
\mathbf{H}_D \mathbf{t}$ with $\mathbf{H}_D$ indefinite, so the zero set
consists of the two asymptote directions $\mathbf{t}$ satisfying
$\mathbf{t}^\top \mathbf{H}_D \mathbf{t} = 0$. In the eigenbasis of
$\mathbf{H}_D$ with $\lambda_1 < 0 < \lambda_2$ these lie at angles
$\pm\arctan\sqrt{-\lambda_1/\lambda_2}$ from the $\lambda_2$
eigenvector. A corner rule can therefore be computed onboard from
$\hat{\mathbf{H}}_D$ alone: inside the gradient-floor ball, evaluate
both estimated asymptotes, discard the one parallel to the incoming
tangent, and exit along the turning branch, sign-chosen by the flow.
Selecting the turning branch at every corner closes the circuit around
one diamond; always selecting the continuing branch reproduces the
published straight-through behavior, so the rule is a second
one-parameter behavior switch, mirroring the park-versus-traverse knob
of the separatrix tracker.

\textbf{Known risk.} The $\hat{\mathbf{H}}_D$ estimate carries a
formation-radius-independent structural bias (third-derivative terms
absent from the quadratic model; see the estimator accuracy results of
the second paper). Eigen-directions survive that bias to within a few
tenths of a degree on the separatrix, but the accuracy of the estimated
asymptote directions at the corners specifically is untested, and near
the corners the eigenvalues themselves are strongly biased. This is the
first thing to measure.

\textbf{Validation plan when revisited.} Implement as a new primitive
(working name \textit{logic\_g\_fable\_step}) alongside the untouched
\textit{logic\_g\_newton\_contour\_pentagon}, following the pattern of
\textit{separatrix\_logic\_fable\_step} (bit-exact equivalence to the
parent law when the knob is off, then the new branch enabled).
Head-to-head at all eight diamond corners plus the domain-center
crossing: success is loop closure around one diamond; metrics are
circumnavigation time and mean $|D|$ per circuit, against the
straight-through baseline of \textit{ow\_clean\_runs.py}.
